\documentclass[11pt,a4paper]{article}
\usepackage[utf8]{inputenc}
\usepackage[T1]{fontenc}
\usepackage[brazil]{babel}
\usepackage[top=1.8cm, bottom=1.8cm, left=2.0cm, right=2.0cm]{geometry}
\usepackage{lmodern}
\usepackage{array}
\usepackage{tabularx}
\usepackage{booktabs}
\usepackage{colortbl}
\usepackage{xcolor}
\usepackage{enumitem}
\usepackage{amsmath}
\usepackage{titlesec}
\usepackage{fancyhdr}

\definecolor{uefsgray}{HTML}{EDEDED}
\definecolor{uefslight}{HTML}{F7F7F7}
\definecolor{uefsdark}{HTML}{1A1A1A}

\titleformat{\section}{\normalfont\bfseries\large}{}{0em}{\MakeUppercase}[\titlerule]
\titlespacing*{\section}{0pt}{12pt}{6pt}

\pagestyle{fancy}
\fancyhf{}
\renewcommand{\headrulewidth}{0pt}
\fancyfoot[L]{\footnotesize UEFS $\cdot$ DCHF $\cdot$ CHF662 Análise Ambiental}
\fancyfoot[R]{\footnotesize Página \thepage}

\setlength{\parindent}{0pt}
\setlength{\parskip}{6pt}

\begin{document}

\begin{center}
  {\bfseries\large UNIVERSIDADE ESTADUAL DE FEIRA DE SANTANA}\\[2pt]
  {\bfseries\normalsize DEPARTAMENTO DE CIÊNCIAS HUMANAS E FILOSOFIA}\\[2pt]
  {\bfseries\normalsize CURSO DE GEOGRAFIA}\\[8pt]
  \rule{\textwidth}{1.2pt}\\[6pt]
  {\bfseries\large ATIVIDADE DE ESTUDO (AULA 11.1)}
\end{center}

\vspace{2pt}

\begin{center}
\renewcommand{\arraystretch}{1.25}
\begin{tabularx}{\textwidth}{|>{\bfseries}p{3.8cm}|X|>{\bfseries}p{2.9cm}|X|}
\hline
\rowcolor{uefsgray}
Componente Curricular & CHF662 Análise Ambiental (60h/a) & Período Letivo & 2026/2 \\
\hline
\rowcolor{uefsgray}
Docente & Prof. Dr. Luiz Diego Vidal Santos & Referência & Aula 11.1 e Capítulo 10 \\
\hline
\end{tabularx}
\end{center}

\vspace{2pt}

\begin{tabularx}{\textwidth}{|X|p{3.2cm}|p{3.5cm}|}
\hline
\textbf{Discente:} & \textbf{Turma:} & \textbf{Data:} \quad / \quad / \quad \\
\hline
\end{tabularx}

\section{Onde estudar}

\noindent\fcolorbox{uefsdark}{uefslight}{%
\begin{minipage}{\dimexpr\textwidth-2\fboxsep-2\fboxrule\relax}
\small
Baseie suas respostas nos slides da \textbf{Aula 11.1 (Indicadores e Monitoramento Ambiental)} e no \textbf{Capítulo 10} da apostila (páginas \textbf{241 a 263}). Consulte as seções sobre dados, indicadores, índices e pirâmide da informação (p.\,242), critérios de seleção (p.\,243), modelos PER e PEIR (p.\,245--246) e ficha do indicador (p.\,247), além dos exercícios de fixação 1 a 5, questões 11 e 12 e itens 1, 3 e 4 da autoavaliação. Atividade sem cálculos numéricos; conclua antes de consultar o gabarito.
\end{minipage}%
}

\section{Questões}

\begin{enumerate}[leftmargin=*, itemsep=6pt]
  \item Diferencie dado bruto, indicador ambiental e índice ambiental, apresentando um exemplo aplicado de cada nível da pirâmide da informação.

  \item Explique por que a utilidade de um diagnóstico ambiental depende da organização dos dados em indicadores vinculados a perguntas de pesquisa e metas de gestão, independentemente do volume bruto de registros coletados.

  \item Apresente os quatro critérios de seleção de indicadores ambientais estudados no capítulo e explique a importância operacional de dois deles para o monitoramento contínuo.

  \item Distinga a estrutura analítica dos modelos PER e PEIR e explique a função da categoria Impacto na separação entre a condição do meio e suas consequências ecológicas ou socioambientais.

  \item Classifique cada situação listada na matriz abaixo segundo as categorias Pressão (P), Estado (E), Impacto (I) ou Resposta (R), registrando uma justificativa sucinta.

  \vspace{4pt}
  \renewcommand{\arraystretch}{1.35}
  \begin{tabularx}{\linewidth}{|p{5.5cm}|c|X|}
  \hline
  \rowcolor{uefsgray}
  \textbf{Situação observada} & \textbf{P, E, I ou R} & \textbf{Justificativa técnica} \\
  \hline
  Abertura de estradas de acesso & & \\
  \hline
  Percentual de solo exposto no talude & & \\
  \hline
  Formação de ravinas e assoreamento a jusante & & \\
  \hline
  Instalação de drenagem e revegetação & & \\
  \hline
  \end{tabularx}
  \vspace{4pt}

  \item Explique por que a fonte, a unidade de medida, o período de coleta, o método e a área de referência precisam acompanhar o registro sistemático de um indicador.

  \item Justifique a exigência de procedimentos padronizados e campanhas comparáveis em programas de monitoramento ambiental e explique qual distorção analítica ocorre quando um trecho sem vistoria recebe registro igual a zero.

  \item Diferencie indicador de execução e indicador de resultado em uma medida de recuperação ambiental, apresentando um exemplo verificável de cada tipo.

  \item Assinale a alternativa correta,
  \begin{enumerate}[label=(\Alph*), leftmargin=1.5em, itemsep=2pt, topsep=3pt]
    \item Um indicador pode ser escolhido sem relação com uma pergunta ambiental, desde que seja fácil de medir.
    \item A sensibilidade é a capacidade de detectar mudanças ambientais e responder às medidas de manejo em prazo útil.
    \item Um índice composto torna desnecessário apresentar os indicadores que o formam.
    \item A mensurabilidade exige equipamentos de alto custo em todas as campanhas.
    \item A relevância de um indicador depende apenas da quantidade de dados produzidos.
  \end{enumerate}

  \item Em uma matriz PEIR, indique a qual componente correspondem a implantação de um programa de monitoramento da fauna e a alteração do traçado de uma linha de transmissão, justificando a classificação.

  \item Explique por que uma baixa taxa agregada de supressão vegetal pode coexistir com fragmentação de habitats e concentração do escoamento superficial em setores críticos da paisagem.

  \item Selecione um problema ambiental recorrente no Semiárido baiano e formule uma pergunta de monitoramento, um indicador compatível com essa pergunta e uma medida de resposta voltada à gestão territorial.
\end{enumerate}

\end{document}
